\documentclass[openright, oneside, a4paper, article, 10pt, brazil]{abntex2}
\usepackage{form-sem-logo}
% --------- %
% Metadados %
% --------- %
\hypersetup{%
	colorlinks=false,
	allbordercolors=white,
	pdftitle = {Declaração institucional para garantia do amplo e livre exercício profissional por parte dos profissionais de Psicologia e Responsáveis Técnicas/os de Psicologia, de acordo com o Código de Ética Profissional da/o Psicóloga/o e legislação cabível},
	pdfauthor = {Conselho Regional de Psicologia de São Paulo},
	pdfsubject = {Amplo e livre exercício profissional},
}
\title{Declaração institucional para garantia do amplo e livre exercício profissional por parte dos profissionais de Psicologia e Responsáveis Técnicas/os de Psicologia, de acordo com o Código de Ética Profissional da/o Psicóloga/o e legislação cabível}
\author{Conselho Regional de Psicologia de São Paulo}
% --------- %
% DOCUMENTO %
% --------- %
\begin{document}
\pagestyle{crp}
\centering
\Form
\setasuspacing{\DoubleSpacing}\centering

\preenctext{14}{Nome_PJ}{40em}{Nome da pessoa jurídica}
\setasuspacing{\OnehalfSpacing}

\section{Declaração institucional para garantia do amplo e livre exercício profissional por parte dos profissionais de Psicologia e Responsáveis Técnicas/os de Psicologia, de acordo com o Código de Ética Profissional da/o Psicóloga/o e legislação cabível}%
\setasuspacing{\DoubleSpacing}\justifying
\vspace*{0.5\baselineskip}

DECLARO, perante o Conselho Regional de Psicologia:
\vspace*{\baselineskip}

\preenctext{14}{nome}{40.3em}{Eu, } \hfill,\\
nascida/o em \dataabrev{14}{1} , \preenctext{14}{RG}{9em}{RG nº } , \preenctext{14}{CPF}{9em}{CPF nº } ,\\
\preenctext{14}{endereco1}{35em}{domiciliada/o à \hfill} \\
\preenctext{14}{endereco2}{\fimdelinha}{\null} ,\\
\preenctext{14}{socia_de}{14em}{sócia/o, proprietária/o e/ou responsável legal pela instituição: \hfill}\\
\preenctext{14}{socia_de-2}{\fimdelinha}{\null} ,\\
\preenctext{14}{cnpj}{9em}{CNPJ nº } , \preenctext{14}{end_institucional}{15em}{com endereço institucional à }\\
\preenctext{14}{end_institucional-2}{\fimdelinha}{\null} ,\\
pleiteante a cadastro/registro de pessoa jurídica de Psicologia neste Conselho Regional de Psicologia, declaro para todos os fins legais e cabíveis:
\vspace{\baselineskip}
%\raggedright

\setasuspacing{\OnehalfSpacing}
\begin{minipage}[c]{\boxlarg}
\CheckBox[%
    name=ciente_resolucao,
    bordercolor=black,
    height=12pt,
    width=12pt]
    {\null}
\end{minipage}%
\begin{minipage}[c]{\linewidth-1.72\boxlarg}
    estar ciente da Resolução nº 16/2019 , que dispõe sobre o registro e cadastro de Pessoas Jurídicas;
\end{minipage}\vspace{\baselineskip}

\begin{minipage}[c]{\boxlarg}
\CheckBox[%
    name=liberdade,
    bordercolor=black,
    height=12pt,
    width=12pt]
    {\null}
\end{minipage}%
\begin{minipage}[c]{\linewidth-1.72\boxlarg}
conceder ampla liberdade na utilização dos métodos e técnicas psicológicas, respeitando sua autonomia profissional,
    os princípios estabelecidos no Código de Ética Profissional do Psicólogo, bem como as demais normativas em
    vigência pertinentes à profissão de Psicologia;
\end{minipage}\vspace{\baselineskip}

\begin{minipage}[c]{\boxlarg}
\CheckBox[%
    name=ciencia_liberdade,
    bordercolor=black,
    height=12pt,
    width=12pt]
    {\null}
\end{minipage}%
\begin{minipage}[c]{\linewidth-1.72\boxlarg}
ter ciência de que o não cumprimento do respeito à liberdade técnica, fundada na ética, dos profissionais que
    trabalham na instituição ao qual sou responsável implicará em sanções administrativas, cíveis e criminais por
    parte dos devidos órgãos competentes;
\end{minipage}\vspace{\baselineskip}

    \begin{minipage}[c]{\boxlarg}
\CheckBox[%
    name=afixar_local,
    bordercolor=black,
    height=12pt,
    width=12pt]
    {\null}
\end{minipage}%
\begin{minipage}[c]{\linewidth-1.72\boxlarg}
afixar em local visível ao público o certificado de registro da inscrição da Pessoa Jurídica, bem como solicitar
    a sua renovação após seu vencimento;
\end{minipage}\vspace{\baselineskip}

    \begin{minipage}[c]{\boxlarg}
\CheckBox[%
    name=documento_comprobatorio,
    bordercolor=black,
    height=12pt,
    width=12pt]
    {\null}
\end{minipage}%
\begin{minipage}[c]{\linewidth-1.72\boxlarg}
encaminhar documento comprobatório ao Conselho Regional de Psicologia de qualquer alteração de seus
atos constitutivos;
\end{minipage}\vspace{\baselineskip}

\begin{minipage}[c]{\boxlarg}
\CheckBox[%
    name=contatos,
    bordercolor=black,
    height=12pt,
    width=12pt]
    {\null}
\end{minipage}%
\begin{minipage}[c]{\linewidth-1.72\boxlarg}
manter os contatos junto ao Conselho Regional de Psicologia atualizados;
\end{minipage}\vspace{\baselineskip}

\begin{minipage}[c]{\boxlarg}
\CheckBox[%
    name=inspecao,
    bordercolor=black,
    height=12pt,
    width=12pt]
    {\null}
\end{minipage}%
\begin{minipage}[c]{\linewidth-1.72\boxlarg}
receber o Conselho Regional de Psicologia durante inspeção de pessoa jurídica, bem como atender às orientações
realizadas pela autarquia;
\end{minipage}\vspace{\baselineskip}

\begin{minipage}[c]{\boxlarg}
\CheckBox[%
    name=publicidade,
    bordercolor=black,
    height=12pt,
    width=12pt]
    {\null}
\end{minipage}%
\begin{minipage}[c]{\linewidth-1.72\boxlarg}
veicular publicidade enquanto pessoa jurídica contendo o número de inscrição no Conselho Regional de Psicologia
e em conformidade com demais normativas vigentes;
\end{minipage}\vspace{2\baselineskip}

\noindent Sem mais,
\vspace{3\baselineskip}

\dataextenso{1}
\vspace{6\baselineskip}

\rule{24em}{0.4pt}

\small Assinatura da/o responsável legal pela instituicão
\vfill

\raggedleft\scriptsize
\textit{Fundamento desta autorização nas seguintes normativas: Lei~nº~4119/1962, Decreto~nº~53.464/1964, Lei~nº~5766/1971,
Decreto~nº~79.822/1977, Resolução~CFP~nº~10/2005 (Código de Ética Profissional do Psicólogo), Resolução~CFP~nº~16/2019}.

\textit{Observação: este documento deverá ser apresentado em papel timbrado da instituição pública ou privada}.
\end{document}